\section{Introduction}
Agent Skills package instructions and helper code so that an agent can perform a reusable task. A reporting Skill, for example, may read a table, calculate totals, write a report and send it to a named recipient. This is useful because the user need not describe each tool call. It also creates a security problem. A malicious package can insert an unrelated credential read, mix private data into a legitimate message, or make the useful work depend on a forbidden action. Studies of malicious Skills and Skill-file injection show why package content should not receive the same trust as the user's request~\cite{liu,skillinject}.

An external gateway can block an action that violates a trusted policy. Blocking is necessary, but it does not by itself finish the user's task. If the gateway stops a credential read in the middle of report generation, the report may remain incomplete. Simply removing that read can also fail: later code may use the missing value, or a condition may make the final send depend on it. The malicious code has created a dependency, but that dependency does not give it permission to read the credential.

Recovery is harder once some useful work has already happened. Suppose the report was written before the block. Replaying the whole Skill can repeat the write. Suppose the report was sent but its acknowledgement was lost. Sending again can produce a duplicate; assuming success can leave the task incomplete when the first send did not happen. A recovery system therefore needs to reason about the current state, not only about the text of the failed script. It also needs to preserve evidence of previous damage. A later correct report cannot turn an earlier disclosure into a safe execution.

An earlier version of our project studied task-scoped checks and legal recovery from a trusted catalogue of replacement scripts. That study established a useful distinction between the original implementation and the user's authorized goal. However, it also supplied an answer to much of the recovery problem: the catalogue already contained a suitable legal route. Its recovery results could not show that a system could construct such a route. This paper removes prepared replacement programs from the main experiment. Trusted task requirements and permissions remain inputs, but the continuation must be built from ordinary operations.

Our question is: \emph{after a Skill is blocked, can a system construct and validate a continuation that completes the trusted task without expanding permissions, losing useful prior work, or repeating side effects?} We test a narrow and executable version of this question. Tasks operate on varying tables and require a correct artifact, one delivery, preservation of inputs and temporary-file cleanup. Malicious helpers are finite programs. A separate evaluator computes the required result from the input data. This gives explicit task truth without using a model judge or treating an attacker comment as an instruction.

The proposed method first builds a residual effect contract. The contract lists unfinished goals, verified intermediate values, confirmed and uncertain effects, remaining tool-call allowance, and dependencies introduced by forbidden data. A continuation constructor uses it to remove completed work, guard uncertain sends with a permitted status query, add bounded retries for confirmed no-effect failures, and cut unneeded computations. When a source condition depends on a forbidden read, the constructor may reuse its branch bodies as candidate material. This is a proposal, not a proof: the same validator used by every method must still reject an incorrect or unsafe result.

If this first construction fails, a model proposes another finite program. Concrete counterexamples guide a bounded refinement loop. These proposals receive no extra authority. In particular, a model cannot repair a failure by adding a new grant or changing the task's success condition. The full method has three candidate-validation slots, just like the ordinary model methods; its source-based first proposal consumes one of those slots. This design permits a practical benefit through fewer model requests, but does not assume that such a benefit exists.

We compare this method with stopping, deletion, replay, ordinary model repair, fresh planning from the task, and candidate generation with the same verifier. We also include a deliberately strong source-free compiler. It directly translates the trusted declarative task into the primitive language and uses the same state handling. This comparator is important because the task specification may already contain enough structure to make original-code repair unnecessary. A method that beats a weak prompt but loses to this compiler cannot be presented as a generally superior synthesis method.

The evaluation separates four questions. First, can a method complete nontrivial blocked tasks without a supplied replacement? Second, do residual goals, dependency cuts, effect history and counterexample feedback contribute beyond the common gateway? Third, do results persist across held-out task structures, new input rows, composed attacks and partial execution? Fourth, what generation and verification costs accompany the gains? Cases with no permission, no remaining allowance, unresolved effects and past disclosure are retained. They test when the system should stop rather than manufacture success.

The contributions are an executable formulation of this restricted recovery problem, an inspectable continuation constructor, and a matched empirical evaluation with full failure and cost records. We do not claim to invent program repair, information-flow control or counterexample-guided synthesis. Nor do we claim safety for arbitrary Python, real messaging services or an entire Skill marketplace. The value of the study depends on which mechanisms work under the stated assumptions and which comparisons remain negative. The report therefore retains strong controls, implementation corrections and results that limit the proposed method.
